% ==============================================================================
% Section 2: Literature Review
% ==============================================================================

\section{Literature Review}
\label{sec:literature}

This section reviews the relevant literature across four interconnected domains: augmented reality in education, intelligent tutoring systems, gamification in learning, and mobile AR applications for reading support. We conclude with an analysis of the research gap that AR-DeC is designed to address.

% ------------------------------------------------------------------------------
\subsection{Augmented Reality in Education}
\label{subsec:ar_education}

Augmented reality, defined by Azuma \cite{azuma1997} as a technology that overlays virtual information onto the real world in real time while maintaining registration with the physical environment, has attracted substantial research attention in educational contexts over the past two decades. The foundational taxonomy proposed by Milgram and Kishino \cite{milgram1994} situates AR on the reality-virtuality continuum, distinguishing it from fully immersive virtual reality by its preservation of the user's connection to the physical environment. This characteristic makes AR particularly suitable for educational applications where contextual grounding in real-world materials---such as textbooks, laboratory equipment, or physical specimens---is pedagogically desirable.

Several comprehensive reviews have documented the growth and impact of AR in education. Wu et al. \cite{wu2013} conducted an influential review covering the period from 2004 to 2012, identifying that AR applications predominantly targeted science education and demonstrated positive effects on learning motivation and spatial understanding. Bacca et al. \cite{bacca2014} extended this analysis through 2014, finding that AR interventions improved learning achievement in 87\% of the studies reviewed, with the strongest effects observed when AR content was tightly aligned with learning objectives. More recently, Ak\c{c}ay{\i}r and Ak\c{c}ay{\i}r \cite{akccayir2017} synthesised 68 studies in a systematic review, confirming positive learning outcomes while identifying challenges related to usability, technical complexity, and the need for stronger theoretical grounding.

The meta-analytic evidence further supports AR's educational effectiveness. Garz\'{o}n et al. \cite{garzonybarragan2020} reported a medium-to-large overall effect size ($g = 0.68$) for AR interventions across educational levels and subject domains. Al-Dokhny et al. \cite{aldokhny2022} corroborated these findings with a meta-analysis of 42 studies, demonstrating that AR significantly improved student achievement compared to conventional instruction, with moderating effects for subject area, educational level, and intervention duration. However, both reviews noted considerable heterogeneity in study designs and outcome measures, suggesting that the field would benefit from more standardised evaluation frameworks.

Despite these promising findings, several researchers have raised important caveats. Makransky et al. \cite{makransky2019} demonstrated that while immersive technologies can increase presence and engagement, they do not automatically improve learning outcomes and may even introduce extraneous cognitive load if not carefully designed. Fowler \cite{fowler2015} argued that many AR educational applications lack pedagogical grounding, functioning as technological novelties rather than theory-driven instructional tools. These critiques underscore the importance of embedding AR within established learning frameworks, as AR-DeC aims to do through its integration with the ARCS model and Bloom's taxonomy.

% ------------------------------------------------------------------------------
\subsection{Intelligent Tutoring Systems}
\label{subsec:its}

Intelligent tutoring systems represent a mature strand of educational technology research with origins in the early work of Carbonell \cite{carbonell1970} on AI-assisted instruction. The classical ITS architecture, as codified by Nwana \cite{nwana1990} and elaborated by subsequent researchers, comprises four interrelated components: a domain model encoding the subject matter knowledge, a student model tracking individual learner states, a tutoring model determining instructional strategies, and an interface model managing the presentation of content.

The effectiveness of ITS has been rigorously established through meta-analytic evidence. Kulik and Fletcher \cite{kulik2016} analysed 50 controlled evaluations and found that ITS produced learning gains roughly equivalent to those achieved by human tutoring, with effect sizes consistently outperforming conventional classroom instruction. VanLehn \cite{vanlehn2011} provided a nuanced comparison showing that while ITS approached the effectiveness of human tutoring for well-structured domains, the gap narrowed further when ITS incorporated sophisticated student modelling and adaptive feedback. These findings align with Bloom's \cite{bloom1984} seminal observation that one-to-one tutoring produces a two-sigma improvement over group instruction, suggesting that ITS can partially bridge this gap at scale.

A critical component of effective ITS is the student modelling approach. Bayesian Knowledge Tracing (BKT), introduced by Corbett and Anderson \cite{corbett1994}, provides a probabilistic framework for estimating a learner's mastery of individual knowledge components based on their history of correct and incorrect responses. BKT models each knowledge component with four parameters: initial knowledge probability, learning rate, guess probability, and slip probability. The model updates its estimate of mastery after each student interaction using Bayesian inference, enabling the system to make informed decisions about when to advance, review, or remediate. Despite its relative simplicity compared to deep learning approaches, BKT remains widely used due to its interpretability, computational efficiency, and strong empirical validation \cite{du2020}.

The integration of ITS with immersive technologies represents an emerging research direction. Nisiotis and Alboul \cite{nisiotis2021} proposed a framework combining VR learning environments with intelligent tutoring for engineering education, demonstrating improved engagement but noting implementation challenges. Sottilare et al. \cite{sottilare2018} developed the Generalized Intelligent Framework for Tutoring (GIFT), which provides a modular architecture for integrating adaptive tutoring with various delivery platforms. However, mobile AR-based ITS implementations remain scarce, particularly those targeting reading comprehension and vocabulary acquisition.

% ------------------------------------------------------------------------------
\subsection{Gamification in Learning}
\label{subsec:gamification}

Gamification, defined by Deterding et al. \cite{deterding2011} as the use of game design elements in non-game contexts, has emerged as a significant strategy for enhancing learner engagement and motivation in educational settings. The theoretical foundations of gamification in education draw on several motivational frameworks, most notably Self-Determination Theory (SDT) \cite{ryan2000}, which posits that intrinsic motivation is sustained when activities satisfy the basic psychological needs of autonomy, competence, and relatedness. Game elements such as points, badges, leaderboards, and progress indicators can support these needs by providing clear goals, immediate feedback, and a sense of achievement.

Keller's ARCS model \cite{keller2010} provides an alternative motivational lens particularly well-suited to instructional design. The model identifies four conditions necessary for sustained motivation: Attention (capturing and maintaining interest), Relevance (connecting content to learner goals and experiences), Confidence (building positive expectations for success), and Satisfaction (providing rewarding outcomes). Each condition maps naturally to specific gamification mechanics: attention can be captured through novel AR visualisations and unexpected rewards; relevance is established by scanning personally-encountered vocabulary in authentic reading materials; confidence is built through adaptive difficulty and progress tracking; and satisfaction is delivered through immediate point awards, badge unlocks, and streak maintenance.

Empirical evidence on gamification's effectiveness in education presents a generally positive but nuanced picture. A substantial body of research indicates that gamification improves engagement metrics such as time-on-task, frequency of voluntary interactions, and task completion rates \cite{bacca2014}. Studies in language learning contexts have shown that gamified vocabulary applications can increase practice frequency by 30--50\% compared to non-gamified alternatives, with the strongest effects observed when game elements are intrinsically linked to the learning activity rather than superficially appended \cite{lan2015}.

However, the literature also identifies important design considerations. Overly competitive elements such as public leaderboards can undermine motivation for lower-performing students, while excessive extrinsic rewards may crowd out intrinsic motivation over time---a phenomenon known as the overjustification effect \cite{ryan2000, csikszentmihalyi1990}. Effective gamification design therefore requires careful calibration of reward schedules, personalisation of challenge levels, and integration of game mechanics with meaningful learning activities rather than arbitrary performance metrics. AR-DeC addresses these concerns by tying gamification elements directly to learning behaviours (scanning, reading, quiz performance) and using the adaptive ITS to maintain an appropriate level of challenge through personalised content selection.

% ------------------------------------------------------------------------------
\subsection{Mobile AR Applications for Reading}
\label{subsec:mobile_ar_reading}

The convergence of mobile computing power, advanced camera systems, and machine learning has enabled a new category of AR applications that can process and augment printed text in real time. Optical character recognition (OCR) technology, particularly on-device implementations such as Google ML Kit, has reached a level of accuracy and speed that enables practical deployment in mobile educational applications. These systems can process camera frames at rates sufficient for interactive use, extracting text with high fidelity across a range of fonts, sizes, and lighting conditions.

Several mobile AR applications have explored the intersection of text scanning and learning support. Translation applications such as Google Translate's camera mode demonstrate the technical feasibility of real-time text recognition and AR overlay, but they operate as simple lookup tools without adaptive learning components. Educational reading applications have explored vocabulary support through digital glossaries and contextual definitions, but typically rely on manual text input rather than camera-based scanning, thereby disrupting the reading flow that OCR-based approaches can preserve \cite{chen2017}.

Research on mobile AR for language learning has yielded encouraging results. Parmaxi \cite{parmaxi2020} reviewed AR applications for language learning and found that mobile implementations were particularly effective for vocabulary acquisition, with the contextual, situated nature of AR providing a form of elaborative encoding that aids retention. Lan \cite{lan2015} demonstrated that 3D virtual environments could enhance contextual vocabulary learning for English as a Foreign Language (EFL) students, suggesting that AR's ability to situate word meanings within visual contexts may offer similar benefits. Ib\'{a}\~{n}ez and Delgado-Kloos \cite{ibanezmolina2018} found that AR applications for STEM learning were most effective when they provided interactive, manipulable visualisations rather than passive overlays.

However, existing mobile AR reading applications share several limitations. First, they typically lack adaptive intelligence---providing the same content to all users regardless of prior knowledge or learning progress. Second, they rarely incorporate gamification elements to sustain engagement over extended use periods. Third, they tend to operate as standalone tools without integration into a broader learning ecosystem that tracks progress and adapts over time. Fourth, most implementations focus on translation or simple definition lookup rather than providing pedagogically structured explanations aligned with the learner's cognitive level.

% ------------------------------------------------------------------------------
\subsection{Research Gap Analysis}
\label{subsec:research_gap}

The preceding review reveals that while AR, ITS, gamification, and mobile OCR each contribute demonstrably to educational outcomes, they have been studied and implemented largely in isolation. \Cref{tab:gap_analysis} summarises the capabilities of existing systems against the key features that AR-DeC integrates.

\begin{table}[htbp]
\centering
\caption{Research gap analysis: comparison of existing approaches with AR-DeC.}
\label{tab:gap_analysis}
\small
\begin{tabular}{lcccccc}
\toprule
\textbf{System Type} & \textbf{OCR} & \textbf{AR} & \textbf{ITS} & \textbf{Gamification} & \textbf{Mobile} & \textbf{Theory} \\
\midrule
AR education apps & -- & \checkmark & -- & -- & \checkmark & Partial \\
ITS platforms & -- & -- & \checkmark & -- & -- & \checkmark \\
Gamified learning apps & -- & -- & -- & \checkmark & \checkmark & Partial \\
Translation apps (e.g., Google) & \checkmark & \checkmark & -- & -- & \checkmark & -- \\
VR/AR + ITS prototypes & -- & \checkmark & \checkmark & -- & -- & \checkmark \\
\midrule
\textbf{AR-DeC (proposed)} & \checkmark & \checkmark & \checkmark & \checkmark & \checkmark & \checkmark \\
\bottomrule
\end{tabular}
\end{table}

The gap analysis reveals three specific deficiencies in the current literature. First, no existing system combines on-device OCR scanning with adaptive tutoring logic, meaning that camera-based reading tools cannot personalise their responses to individual learner needs. Second, the integration of gamification with AR learning environments remains largely unexplored beyond surface-level point systems, with no implementations grounding game mechanics in established motivational theory such as the ARCS model. Third, mobile AR applications for vocabulary support have not incorporated student modelling techniques such as BKT, which would enable longitudinal tracking of word mastery and adaptive content sequencing.

AR-DeC is designed to address all three gaps simultaneously, providing a unified mobile platform that integrates OCR-based text scanning, BKT-driven adaptive tutoring, AR contextual visualisation, and ARCS-grounded gamification. The following sections detail the theoretical framework, system architecture, and technical implementation of this integrated approach.
